\chapter*{Agradecimientos}
En primer lugar, quiero expresar mi más sincero agradecimiento a mis tutores,
Antonio Ruiz-Cortés y Manuel Resinas. Gracias por la confianza que depositasteis
en mí al abrirme las puertas de este proyecto, y por una guía que me ha hecho
crecer como ingeniero y me ha enseñado a mirar cada problema con más rigor.
Hago extensivo este agradecimiento a Ramón Gavira y a Daniel Ruiz, cuyo trabajo
es el punto de partida de este. Sin las herramientas que construyeron, este
trabajo no habría sido posible.

También quiero dar las gracias a Daniel Romero. Ya me ayudó a sacar adelante mi
Trabajo Fin de Grado, y quién me iba a decir entonces que hoy trabajaríamos
juntos y compartiríamos casa. Has vivido de cerca el día a día de este trabajo,
con sus noches largas y sus dudas, y siempre has tenido una palabra para que
siguiera adelante. Gracias por seguir ahí, dentro y fuera del trabajo.

En un lugar aparte, quiero mencionar a mi pareja, Leticia. Gracias por tu apoyo
y por haber estado conmigo en las etapas más difíciles de este camino, sobre
todo en aquellas en las que yo mismo dudaba de poder terminarlo. Cuando el peso
se hacía más grande, tú lo hacías más ligero. No podría haber cerrado mejor este
trabajo que contigo a mi lado.

Finalmente, dedico el agradecimiento más profundo a mi familia. A mis padres,
por seguir siendo la base sobre la que se sostiene todo lo que hago, y por
enseñarme que las cosas se terminan aunque cuesten. A mi hermano, porque a
pesar de la distancia te siento cerca en cada paso, y cada logro mío sigue
siendo también tuyo. Este trabajo, como el anterior, es en gran parte vuestro,
y saber que estáis orgullosos de mí sigue siendo el mayor motor que tengo.
